\scanpage{12}
\mspage{\handwritten{30}}
% A single long oblique stroke crosses the opening through the theorem heading.
\deleted{\deleted{on a nécessairement} puisque ses composantes connexes sont ouvertes ($S$ étant localement noethérien). Mais alors on a nécessairement $Z=J_S$, où $J$ est une partie de $I$, et on a \deleted{\unsure{$Z=$|$Z=K$}} donc $U=\emptyset$ ou $U=S$, suivant que $J$ est infini ou fini, ce qui donne la conclusion voulue.

\uline{Corollaire} 5.1\deleted{\unsure{3}}\handwritten{4}. Soi\deleted{t}\inserted{ent} $S$ un préschéma localement noethérien et \uline{normal}, $\deleted{\unsure{R}}\handwritten{P}$ un $S$-préschéma \deleted{localement} constant t\deleted{\unsure{r}}o\deleted{\unsure{o}}rdu quasi-isotrivial. Alors les composantes connexes de $\deleted{R}\handwritten{P}$ sont finies sur $S$.

Comme $\deleted{R}\handwritten{P}$ est étale sur $S$, on sait que $\deleted{R}\handwritten{P}$ est normal (SGA I 9.5.) donc ses composantes \deleted{irréductibles} connexes sont ses composantes irréductibles. On peut évidemment supposer $S$ connexe donc irréductible, soit $s$ son point générique. Comme $R$ est plat sur $S$, ses points maximaux, i.e. les points\handwritten{\unsure{'}} génériques $x_i$ des composantes irréductibles $R_i$ de $R$, sont les points maximaux de $R_s$, et la fibre \deleted{de $R$} gén\handwritten{é}rique $R_{i\,s}$ de $R_i$ est la composante irréductible de $R_s$\handwritten{,} adhérence de $x_i$ dans $R_s$. Or $R_s$ étant étale sur $k(s)$, est somme disjointe de spectres d\deleted{\unsure{e}}'extensions finies séparables de $k(s)$, donc \deleted{les} pour tout $i$, $R_{i\,s}$ est fini (en fait, réduit à un seul point). Donc 5.1\deleted{2}\handwritten{3}. s'applique et montre que les $R_i$ sont finis sur $S$.

\uline{Remarque} 5.1\deleted{\unsure{4}}\handwritten{5}. Le raisonnement s'applique, plus généralement, si on suppose seulement $S$ géométriquement unibranche au lieu de normal, i.e. que le morphisme \deleted{\unsure{S}}canonique $\overset{\handwritten{\sim}}{S}\handwritten{\longrightarrow}S$, où $\overset{\handwritten{\sim}}{S}$ est le normalisé de $S_{\text{réd}}$\handwritten{,}\deleted{\unsure{,}} est radiciel. Donc le théorème 5.1\deleted{5}\handwritten{6}. ci-dessous est valable également sous cette condition plus générale.

\uline{Théorème} 5.1\deleted{5}\handwritten{6}. Soit $S$ un préschéma localement noethérien normal. Alors tout groupe de type multiplicatif} \deleted{\unsure{et de type fini|de type fini} \unsure{présentation|quasi-isotrivial}} \inserted{et de type fini} sur $S$\deleted{\unsure{, est et en parti|, et et en partic}} est \uline{isotrivial}.

On peut supposer en effet \inserted{$S$ connexe, donc} $H$ de type constant $M$. Il suffit d'appliquer 5.13 . avec $P=\underline{\operatorname{Isom}}_{S\text{-gr}}(H,G)$, où $G=D_S(M)$\handwritten{,}
\handwritten{en raisonnant comme dans 5.1\unsure{2|1}. (ii) $\Longrightarrow$ (i); \deleted{\illegible{2}} \inserted{On peut aussi} appliquer \unsure{à|ou} \deleted{\unsure{avec}} 5.11. \deleted{\illegible{4}}, en utilisant 5.14. pour $R=D_S(\underline{H})$.}
\margin{\handwritten{5.11. \textcircled{\unsure{M|M$^0$}}.}}
% archive sheet number: 11

